\section{Introduction}

Quantum chromodynamics (QCD) connects hadronic matter to its fundamental 
constituents, quarks, and gluons. Many aspects of QCD are poorly understood, in 
spite of four decades of intense 
experimental and theoretical effort. As part of this effort, deep inelastic 
scattering (DIS) of leptons from nucleons has played a central role in 
establishing QCD as the reigning theory of the strong interaction and continues 
to serve as a mainstay for attempts to unravel QCD's mysteries 
(for complete reviews see, for example, Refs.~\cite{pdg12} and 
\cite{Collins:2011zzd}).

Theoretically, inclusive DIS cross sections are 
separated into leptonic and hadronic tensors, which capture the 
electroweak and strong dynamics of the scattering process, respectively. The 
hadronic tensor factorizes into a convolution of infrared-safe perturbative 
coefficients and parton distribution functions (PDFs), 
which incorporate the low-energy QCD physics and therefore must be determined 
nonperturbatively. PDFs are independent of the scattering process--but 
depend on the target nucleon, while the perturbative coefficients are 
independent of the external scattering states.

PDFs are important for two reasons. First, they furnish 
direct knowledge of the constituents of hadronic states---the dominant form of 
visible matter in the universe---in terms of the fundamental theory of the 
strong force. Second, they provide constraints on hadronic 
backgrounds at collider experiments, such as the Large Hadron Collider. 
These backgrounds affect the sensitivity of a variety of high energy 
experiments, including studies of the properties of the Higgs boson 
\cite{Thorne:2011kq,Alekhin:2010dd,Alekhin:2011ey}, and searches for heavy $W'$ 
and $Z'$ boson production \cite{Brady:2011hb}.

PDFs cannot currently be calculated directly from QCD using \emph{ab 
initio} lattice QCD, because they are defined 
in terms of light-cone matrix elements that are not directly 
accessible on Euclidean lattices. Hence, PDFs
are usually extracted from global analyses of different experiments  
\cite{Jimenez-Delgado:2014xza,Jimenez-Delgado:2013boa,Ball:2013lla,
Arbabifar:2013tma,Owens:2012bv,Accardi:2011fa,Leader:2010rb,Blumlein:2010rn, 
Hirai:2008aj}. 

Naturally a direct nonperturbative computation of PDFs is desirable: with 
sufficient precision, such a calculation would further constrain global fits in 
regions that may be 
experimentally inaccessible.

Until recently, lattice calculations have focused on determining the Mellin 
moments of PDFs, which are related to matrix elements of twist-2
operators that \emph{can} be determined on the lattice 
\cite{Bali:2012av,Braun:2008ur,Guagnelli:2004ga}. The lattice regulator breaks 
Lorentz symmetry, which induces radiative mixing
between operators of different mass dimensions. The resulting power divergences 
in the lattice spacing completely 
obscure the continuum limit. Moments up to the fourth moment 
can 
be extracted by carefully choosing the external momenta and operators, but this 
significantly 
limits the precision with which one can 
extract meaningful results for PDFs \cite{Detmold:2003rq,Detmold:2001dv}. 
Beyond 
the fourth moment, power divergent mixing is inevitable and this method cannot 
provide reliable calculations.

Ji recently proposed a new approach to directly extract PDFs from lattice 
calculations \cite{Ma:2014jla,Ji:2013dva}, 
based on a large-momentum effective theory \cite{Ji:2014gla}. Preliminary 
results first appeared in Ref.~\cite{Lin:2014zya}, but a number of issues 
remain, including a complete understanding of the renormalization of the 
relevant lattice matrix elements and the practical ability to resolve 
sufficiently large momenta on the lattice.

Here we propose a new formalism that removes mixing in the continuum limit and, 
in 
principle, enables the extraction of higher moments of PDFs from lattice QCD. 
We call this formalism the ``smeared operator product expansion'' (sOPE). We 
expand continuum matrix elements in a basis of locally-smeared lattice 
degrees of 
freedom. The resulting matrix elements are functions of two scales, the 
smearing 
scale and the renormalization scale. The sOPE provides the framework necessary 
to relate these matrix elements, via Wilson coefficients, to 
phenomenologically useful quantities, such as the hadronic tensor of DIS.

Smearing has been widely used in lattice QCD to reduce ultraviolet 
fluctuations, partially restore rotational symmetry and thereby systematically 
improve the precision of lattice calculations 
\cite{Morningstar:2003gk,Hasenfratz:2001hp,Bernard:1999kc,Albanese:1987as}. For 
a pedagogical overview of smearing in lattice 
calculations, see, for example, Ref.~\cite{Gattringer:2010zz}. In the sOPE, we 
implement smearing 
via the gradient flow, a classical evolution of the fields in a new dimension, 
the flow time, toward the stationary points of the action 
\cite{Luscher:2013cpa,Luscher:2011bx,Lohmayer:2011si,Luscher:2010iy,
Narayanan:2006rf}. 
Matrix elements determined nonperturbatively on the lattice require no further 
renormalization, up to a fermionic wave function renormalization 
\cite{Luscher:2013cpa}, and remain 
finite, provided the smearing scale is kept fixed in physical units in the 
continuum limit \cite{Makino:2014sta,Luscher:2011bx}. There are two further 
advantages: the gradient flow allows one to use smearing lengths 
of only one or two lattice spacings, much smaller than hadronic length scales 
on typical lattices, and therefore does not distort the low energy physics 
\cite{Davoudi:2012ya}, and also the gradient flow is computationally very 
cheap. 

The gradient flow is now well-established as a tool to study lattice 
gauge theories, with applications from scale-setting, 
\emph{i.e.}~determining the lattice spacing in physical units 
\cite{Borsanyi:2012zr,Borsanyi:2012zs,Bazavov:2013gca,Fodor:2014cpa,
Cheng:2014jba}, to studying renormalization in 
lattice gauge theories. For example, 
the gradient flow has been used to define finite-volume renormalization schemes 
for the strong coupling constant
\cite{Fodor:2012td,Fritzsch:2013yxa,Fritzsch:2013hda, 
Ramos:2013gda,Rantaharju:2013bva}, for operator 
renormalization \cite{Monahan:2013lwa}, and to understand the nonperturbative 
scale dependence of renormalized matrix elements \cite{Luscher:2014kea}. 
Related work has used Ward identities to study chiral fermions on 
the lattice \cite{Shindler:2013bia} and the energy-momentum tensor 
\cite{Suzuki:2013gza,DelDebbio:2013zaa, Makino:2014taa}.

In this paper we introduce the gradient flow for a 
single real scalar field with quartic interactions and outline some of the 
properties of the sOPE applied to real scalar field theory.
Scalar field theories describe some important critical 
phenomena in nature, such as the antiferromagnetic phase transition, and 
have a long history as a testing ground for fundamental ideas in 
quantum field theory in four dimensions. For our purposes, the chief advantage 
is the simplicity of Euclidean scalar field theory, which lays bare the 
structure of the sOPE and highlights the points of similarity and contrast with 
the local operator product expansion (OPE). 

In the next section, we review Wilson's OPE and introduce the sOPE. We then 
apply the sOPE to scalar field theory in Sec.~\ref{sec:phi4}, illustrate how 
the gradient flow removes power-divergent mixing in Sec.~
\ref{sec:lattsmear}, 
and calculate 
the perturbative Wilson coefficients to two loops in Sec.~\ref{sec:wc}. 
Finally, in Sec.~\ref{sec:rgeqns}, we study the scale dependence of the sOPE 
through renormalization group 
equations for the Wilson coefficients.
